%LTeX: language=fr-FR
\documentclass[12pt,a4paper]{article}
\usepackage[utf8]{inputenc}
\usepackage[T1]{fontenc}
\usepackage[french]{babel}
\usepackage[hmargin=1.6cm, vmargin=1.1cm]{geometry}
\usepackage{amsmath,amssymb}
\usepackage{graphicx}
\usepackage{enumitem}
\usepackage{tcolorbox}
\usepackage{xcolor}
\usepackage{multicol}
\usepackage{makecell}
\usepackage{tasks}
\usepackage{tikz}

\tcbuselibrary{skins}

\definecolor{bandeau}{HTML}{E8EEF7}
\definecolor{titrebleu}{HTML}{1A4FA0}

% Bandeau de titre pour chaque exercice
\newtcbox{\exobox}{
  colback=bandeau, colframe=bandeau, boxrule=0pt,
  arc=0pt, left=4pt, right=4pt, top=3pt, bottom=3pt,
  fontupper=\bfseries, width=\linewidth, halign=left
}

% \exotitre{numéro}{titre}{barème} : bandeau seul (utilisable dans une minipage)
\newcommand{\exotitre}[3]{%
  \noindent\exobox{Exercice #1 --- #2 \hfill (#3 points)}\par
  \vspace{0.4em}
}
% \exercice{numéro}{titre}{barème} : espace + bandeau
\newcommand{\exercice}[3]{%
  \vspace{0.35em plus 0.5fill}
  \exotitre{#1}{#2}{#3}
}
% Exercice avec figure : bandeau et questions dans une minipage à gauche, figure à droite alignée en haut
\newcommand{\exoespace}{\vspace{0.35em plus 0.5fill}}

\pagestyle{empty}
\setlist{topsep=2pt, itemsep=2pt, parsep=0pt}
\setlength{\parskip}{3pt}
\setlength{\parindent}{0pt}

\begin{document}

\begin{center}
  {\large\bfseries Contrôle --- Fonctions affines et suites}\\[0.4em]
  {\itshape 1\textsuperscript{ère} ST2S \quad $\bullet$ \quad Durée~: 55 minutes \quad $\bullet$ \quad Calculatrice non autorisée \quad $\bullet$ \quad \textbf{Sujet A}}
\end{center}

\vspace{0.3em}
\noindent\textit{La qualité de la rédaction et la justification des réponses seront prises en compte dans la notation.}

% ================= Automatismes =================
\exercice{1}{Automatismes}{4}
Les trois questions sont indépendantes. Détailler les calculs.
\begin{enumerate}
  \item Calculer $\sqrt{2} \times \sqrt{18}$.
  \item Résoudre l'équation $(x + 3)(2x - 8) = 0$.
  \item Un article coûte $200$~euros. Son prix augmente de $10\,\%$, puis il baisse de $10\,\%$. Calculer le prix final de l'article.
\end{enumerate}

% ================= Fonctions affines =================
\exercice{2}{Fonction affine}{6}
Soit $f$ la fonction définie sur $\mathbb{R}$ par $f(x) = -2x + 8$.
\begin{enumerate}
  \item Donner le coefficient directeur et l'ordonnée à l'origine de la droite qui représente $f$.
  \item Donner le sens de variation de $f$ sur $\mathbb{R}$. Justifier.
  \item Calculer $f(3)$.
  \item Déterminer, par le calcul, l'antécédent de $14$ par $f$.
  \item Dresser le tableau de signes de $f(x)$ sur $\mathbb{R}$.
\end{enumerate}

\exoespace
\noindent\begin{minipage}[t]{0.56\linewidth}
\vspace{0pt}
\exotitre{3}{Droites}{4}
\begin{enumerate}
  \item On donne les points $A(1~;~1)$ et $B(3~;~5)$. Calculer le coefficient directeur de la droite $(AB)$, puis déterminer son équation réduite.
  \item Lire graphiquement l'équation réduite de la droite $d$ représentée ci-contre.
\end{enumerate}
\end{minipage}
\hfill
\begin{minipage}[t]{0.41\linewidth}
\vspace{0pt}
\centering
\begin{tikzpicture}[scale=0.5]
  \draw[gray!40, very thin, step=1] (-2,-2) grid (7,5);
  \draw[->, thick] (-2,0) -- (7.4,0) node[below right] {\small $x$};
  \draw[->, thick] (0,-2) -- (0,5.4) node[above left] {\small $y$};
  \foreach \x in {-2,-1,1,2,...,7} \node[below, font=\scriptsize] at (\x,0) {$\x$};
  \foreach \y in {-2,-1,1,2,...,5} \node[left, font=\scriptsize] at (0,\y) {$\y$};
  \node[below left, font=\scriptsize] at (0,0) {$0$};
  \draw[very thick, blue!70!black, domain=-2:7] plot (\x,{-0.5*\x+3});
  \node[above right] at (1,2.5) {$d$};
\end{tikzpicture}
\end{minipage}

% ================= Suites =================
\exercice{4}{Une suite en contexte}{3}
Un EHPAD accueille $60$ résidents en 2026. Chaque année, il accueille $4$ résidents de plus. On note $r_n$ le nombre de résidents en $2026 + n$. Ainsi $r_0 = 60$.
\begin{enumerate}
  \item Calculer $r_1$ et $r_2$.
  \item Compléter la relation de récurrence~: $r_{n+1} = r_n + \dots$
  \item On admet que $r_n = 60 + 4n$. Calculer $r_{10}$, puis interpréter ce résultat.
\end{enumerate}

\exoespace
\noindent\begin{minipage}[t]{0.56\linewidth}
\vspace{0pt}
\exotitre{5}{Termes et nuage de points}{3}
\begin{enumerate}
  \item La suite $v$ est définie par $v_0 = 3$ et $v_{n+1} = 2 \times v_n - 1$. Calculer $v_1$, $v_2$ et $v_3$ en détaillant les calculs.
  \item On a représenté ci-contre les premiers termes d'une suite $w$. Lire $w_0$ et $w_3$, puis conjecturer le sens de variation de la suite $w$.
\end{enumerate}
\end{minipage}
\hfill
\begin{minipage}[t]{0.41\linewidth}
\vspace{0pt}
\centering
\begin{tikzpicture}[x=0.75cm, y=0.29cm]
  \draw[gray!40, very thin, xstep=1, ystep=1] (0,0) grid (6,13);
  \draw[->, thick] (0,0) -- (6.5,0) node[below right] {\small $n$};
  \draw[->, thick] (0,0) -- (0,13.6);
  \foreach \x in {1,...,6} \node[below, font=\scriptsize] at (\x,0) {$\x$};
  \foreach \y in {2,4,...,12} \node[left, font=\scriptsize] at (0,\y) {$\y$};
  \node[below left, font=\scriptsize] at (0,0) {$0$};
  \foreach \n/\w in {0/12, 1/9, 2/7, 3/6, 4/5, 5/4.5} \fill[red] (\n,\w) circle (3pt);
\end{tikzpicture}
\end{minipage}

\end{document}
